\documentclass[UTF8,fontset=none,11pt,a4paper]{ctexart}
\usepackage{ipara-handbook}
\usepackage{graphicx}
\handbooksetup{DS-09 查找与有序索引}{408 · 数据结构 DS-09}{Order · Search · Balance · Update}
\begin{document}
\handbooktitle{查找与有序索引}{维护多少秩序，才能换取多少查询能力}{408 · 数据结构 Topic DS-09}

\begin{corebox}{母问题：查询加速从哪里来，又由谁付费？}
有序索引的生成链是
\[
\text{Order/Index}\to\text{Search Reduction}\to\text{Update Repair}\to\text{Height Invariant}\to\text{Cost Tradeoff}.
\]
顺序查找几乎不预处理；折半查找要求随机访问和有序区间；BST 把次序放进树路径；平衡树用旋转或分裂把高度约束在对数级。查询越强，通常越要在插入、删除、空间或维护时间上付费。
\end{corebox}
\begin{methodbox}{本册定位与 Owner 边界}
\begin{itemize}
\item \kw{Owns：}顺序/分块/折半查找、BST、AVL/红黑树的有序性与高度维护、B 树的查找与插删、B+ 树的查找结构。
\item \kw{Uses：}Atlas Foundation 成本向量；DS-B02 负责按 workload 比较有序索引与 Hash。
\item \kw{Stops：}Hash 映射归 DS10；磁盘块 I/O 的外部排序归 DS12；具体工程库实现只作 Extension。
\end{itemize}
\end{methodbox}
\tableofcontents\newpage

\section{查找契约先固定}
“查找”要先回答返回的是存在性、位置、前驱、首个重复 key 还是范围边界。折半查找依赖一个明确的区间不变量，例如目标若存在必在 $[low,high]$ 中；每轮比较后必须缩小区间，空区间才是失败证明。重复 key 若不固定取左界/右界，算法结果虽都可能“找到”却不满足接口。

\subsection{顺序、分块与折半：三种秩序预算}
顺序查找不要求数据有序，逐项排除候选。分块查找把表分成若干块，维护“块间有序、块内可无序”的两级结构：先在索引表定位可能的块，再在块内顺序查找。其平均成本可写成
\[
ASL\approx ASL_{index}+ASL_{block},
\]
因此块数增加会让索引阶段变长、块内阶段变短；不能只优化其中一项。折半查找则要求整体有序且支持随机访问，每轮用中点把候选区间近似减半。

三者逐步增加维护秩序以缩小查询候选：无索引、稀疏块索引、全局有序随机访问。若数据频繁插入，维持整体连续有序的移动成本可能抵消查询收益。

\subsection{ASL 与判定树：把“快”还原成比较次数}
查找表是逻辑集合，不等于某一种存储结构；静态表可预先排序和建索引，动态表则要把更新维护写入成本。若成功查找记录 $i$ 的关键字比较次数为 $C_i$、访问概率为 $P_i$，平均查找长度为
\[
ASL_{succ}=\sum_i P_iC_i.
\]
等概率时就是比较次数的算术平均。失败查找的平均不能机械除以 $n$：基于比较的有序查找由 $n$ 个关键字切出 $n+1$ 个失败区间，判定树有 $n$ 个内部节点和 $n+1$ 个外部节点；应按实际失败入口数加权。散列查找则由哈希函数值域决定失败入口，不能把数组容量直接当分母。

判定树把每次比较画成内部节点，把空区间或空指针画成失败外部节点。成功比较次数是根到内部节点的层数，失败比较次数是到外部节点路径上的内部节点数。这个模型同时解释了顺序查找的单支链、折半查找的近似平衡树，也能验证“复杂度只数关键字比较，不含地址计算”的题目口径。

\subsection{顺序查找的三种边界}
无序表逐项扫描；逆向扫描可把下标 $0$ 留作失败返回位置。监视哨把目标临时写在哨兵位置，省掉每轮越界判断，但关键字比较次数的渐近量级不变，失败时还会多比较一次。若表有序，扫描到第一个大于目标的 key 即可失败；此时成功 $ASL=(n+1)/2$，失败区间的比较次数为 $1,2,\ldots,n,n$，而不是把最后一个区间误算成 $n+1$。

\subsection{折半查找：区间不变量与左右边界}
标准模板维护闭区间 $[low,high]$：循环条件是 $low\le high$，中点可写为 $low+(high-low)/2$，排除中点后必须使用 $mid-1$ 或 $mid+1$，否则区间可能不缩小。它同时要求全局有序和 $O(1)$ 随机访问；链表即使有序，找中点也要线性走指针。

若重复 key 的接口要求“第一个位置”，命中后不能立即返回，而应记录答案并令 $high=mid-1$；要求“最后一个位置”则令 $low=mid+1$。更一般地，先找第一个满足 $a[i]\ge key$ 的位置（lower bound），或第一个满足 $a[i]>key$ 的位置（upper bound），再用边界区间判断是否真的命中。任何二分题都应先写谓词的单调性，再决定收缩方向。

折半查找的判定树根到内部节点的层数就是成功比较次数；失败查找落到 $n+1$ 个外部叶，路径上的内部节点数就是失败比较次数。完全二叉判定树高度约为 $\lceil\log_2(n+1)\rceil$，但 $n$ 不满一整层时各外部路径长度不同，失败 ASL 不能只写一个高度。给定具体表，直接把根、中点、左右子区间画成树再求和，能避免把“成功 ASL”和“失败 ASL”共用分母。

`mid=(low+high)/2` 与 `mid=(low+high+1)/2` 分别偏向下界或上界；当循环不变式是“寻找最后一个满足谓词的位置”时，必须选择与更新式匹配的中点，否则两元素区间可能在 `low=mid` 时不再缩小。代码层面可写成 `low + (high-low)/2` 防止整数加法溢出。

\subsection{答案二分：搜索的是单调可行域}
有些题不直接查数组位置，而是查最小容量、最短时间或最大允许阈值。先构造判定函数 $P(x)$，证明它在候选答案区间上单调，例如“容量 $x$ 能否在限定天数内完成运输”：若 $x$ 可行，则更大容量也可行。此时目标是找第一个可行值，二分模板与 lower bound 相同；若目标是最后一个可行值，则使用相反边界和上取整中点。

完整复杂度为
\[
O\bigl(C_{check}\log(U-L+1)\bigr),
\]
其中 $[L,U]$ 必须覆盖答案，$C_{check}$ 是一次可行性检查成本。上下界不是随便取：装载容量的下界至少是最大单件重量、上界可取总重量；若 $P$ 并不单调，二分即使代码终止也没有正确性。整数乘加可能溢出时，判定函数应采用扩展类型或提前比较，不能只修正中点公式。

\subsection{分块查找的两级不变量}
长度为 $n$ 的表分成 $b$ 个块，块内可无序但块间必须整体有序；索引项通常保存该块最大关键字和起始地址，先找第一个 $max_i\ge key$ 的块，再在该块内顺序扫描。块内无序换来了插入删除不必维持全表移动，索引查找可用顺序或折半。
若每块 $s$ 个元素且索引、块内都用顺序查找，成功查找平均成本近似
\[
ASL=\frac{b+1}{2}+\frac{s+1}{2},\qquad bs=n.
\]
由均值不等式可知两段工作量平衡时最优，即 $b=s=\sqrt n$；索引改为折半后，最优块长会改变，不能沿用这个结论。

\section{BST：次序下沉到路径}
BST 对每个节点保持左子树 key 小于节点、右子树 key 大于节点（重复值策略需显式约定）。查找、插入、删除都沿一条根叶路径；复杂度是 $O(h)$，不是无条件 $O(\log n)$。退化链的高度为 $n$，所以平衡条件不是装饰而是复杂度证明的前提。
\begin{examplebox}{删除的三种结构分支}
叶节点直接删除；单孩子用孩子替代节点；双孩子从中序后继/前驱取替代 key，再在其原位置删除。每一步都要恢复局部 BST 有序性，不能只交换值而忘记释放或接回子树。
\end{examplebox}

BST 的中序遍历必须得到非降 key 序列，这是验证顺序不变量的最短证据；若重复 key 采用“相等统一放右侧”等规则，验证时也必须使用同一严格/非严格关系。求最近邻、前驱或后继不必遍历整棵树：沿搜索路径维护当前最佳候选即可。构造树时，先序/后序给根，中序把左右子树切开；若 key 重复，序列通常不能唯一确定一棵树，题面必须给出额外规则。

\subsection{BST 上的三类派生查询：秩、累加与全局边界}
BST 的有序性不只用于查找 key，还能把一次中序/反中序扫描变成派生状态：
\begin{itemize}
\item \kw{第 $k$ 小：}中序访问天然按非降序产生 key，维护计数器，计数到 $k$ 即可停止；若树高为 $h$，沿路径并访问前 $k$ 个节点的成本可写为 $O(h+k)$，不提前停止则为 $O(n)$。重复 key 的“第 $k$ 个”必须先约定按节点还是按不同 key 计数。
\item \kw{大于等于当前值的累加树：}反中序（右—根—左）按降序访问，维护累计和并把当前 key 替换为“原值加上已见更大值”；必须先保存原 key，不能在比较结构仍需使用时覆盖后再递归。
\item \kw{验证是否为 BST：}不能只检查每个节点与直接孩子的局部大小；应沿递归传递严格的 $(lower,upper)$ 开区间，或检查中序序列是否满足题目约定的严格/非严格单调性。局部孩子比较无法发现“孙节点越过祖先边界”的反例。
\end{itemize}
这些应用的共同前提是树结构确实满足 BST 不变量。若输入只是任意二叉树，先做验证或改用显式排序；不能因为“写了中序”就自动拥有有序语义。

\section{AVL 与红黑树：都控制高度，但修复契约不同}
AVL 约束任一节点左右子树高度差不超过 1。插入后从修改点向根寻找第一个失衡祖先，根据较高子树方向形成 LL/RR/LR/RL；旋转后必须同时恢复 BST 中序有序和高度约束。

四种插入失衡可压缩为两次选择：先看失衡节点的重子树方向，再看重子树根的重子树方向。LL 与 RR 各一次单旋，LR 与 RL 先对孩子反向旋转再对失衡节点旋转。删除可能让多层祖先连续失衡，不能套用“插入只修第一个节点就停止”的经验；每次旋转后都要重新计算高度。

\begin{center}
\resizebox{0.94\linewidth}{!}{\providecolor{themebg}{HTML}{FAFAF7}
\providecolor{themefg}{HTML}{111111}
\providecolor{themegray}{HTML}{666666}
\providecolor{themecurve}{HTML}{1D63B8}
\providecolor{themealert}{HTML}{C53030}
\providecolor{themeamber}{HTML}{B86E00}
\begin{tikzpicture}[font=\scriptsize,color=themefg,text=themefg,>=Stealth]
\tikzset{
  n/.style={circle,draw=themefg,fill=themebg,minimum size=6.8mm,inner sep=0pt,line width=.8pt,font=\small},
  hot/.style={n,draw=themealert,fill=themealert!10!themebg},
  edge/.style={themegray,line width=.85pt},
  rot/.style={->,themecurve,line width=1.1pt},
  sub/.style={draw=themegray,rounded corners=2pt,minimum width=8mm,minimum height=5mm,inner sep=1pt,text=themegray}
}
\node[anchor=west,font=\bfseries\large] at (0,10.9) {AVL 四类旋转：只改局部父子关系，中序次序始终保持};
\node[anchor=west,font=\small,text=themegray] at (0,10.35)
  {四个面板只改变“新增路径的两次方向”。旋转后统一得到 $T_1<B<T_2<C<T_3<A<T_4$ 或其对称的同一中序次序。};

% helper macro-like manual panels
% LL
\node[font=\bfseries,text=themeamber] at (3.7,9.45) {LL：对 $A$ 右旋};
\node[hot] (lla) at (1.8,8.5) {$A$}; \node[n] (llb) at (1.0,7.55) {$B$}; \node[n] (llc) at (.45,6.55) {$C$};
\node[sub] (llt1) at (.05,5.7) {$T_1$}; \node[sub] (llt2) at (.85,5.7) {$T_2$}; \node[sub] (llt3) at (1.55,6.55) {$T_3$}; \node[sub] (llt4) at (2.65,7.55) {$T_4$};
\draw[edge] (lla)--(llb); \draw[edge] (lla)--(llt4); \draw[edge] (llb)--(llc); \draw[edge] (llb)--(llt3); \draw[edge] (llc)--(llt1); \draw[edge] (llc)--(llt2);
\draw[rot] (3.0,7.35)--node[above,font=\scriptsize]{右旋}(4.0,7.35);
\node[n,draw=themecurve] (ll2b) at (5.2,8.5) {$B$}; \node[n] (ll2c) at (4.35,7.4) {$C$}; \node[n] (ll2a) at (6.05,7.4) {$A$};
\node[sub] (ll2t1) at (3.85,6.35) {$T_1$}; \node[sub] (ll2t2) at (4.85,6.35) {$T_2$}; \node[sub] (ll2t3) at (5.55,6.35) {$T_3$}; \node[sub] (ll2t4) at (6.55,6.35) {$T_4$};
\draw[edge] (ll2b)--(ll2c); \draw[edge] (ll2b)--(ll2a); \draw[edge] (ll2c)--(ll2t1); \draw[edge] (ll2c)--(ll2t2); \draw[edge] (ll2a)--(ll2t3); \draw[edge] (ll2a)--(ll2t4);

% RR
\node[font=\bfseries,text=themeamber] at (11.8,9.45) {RR：对 $A$ 左旋};
\node[hot] (rra) at (9.0,8.5) {$A$}; \node[n] (rrb) at (9.8,7.55) {$B$}; \node[n] (rrc) at (10.35,6.55) {$C$};
\node[sub] (rrt1) at (8.15,7.55) {$T_1$}; \node[sub] (rrt2) at (9.25,6.55) {$T_2$}; \node[sub] (rrt3) at (9.95,5.7) {$T_3$}; \node[sub] (rrt4) at (10.75,5.7) {$T_4$};
\draw[edge] (rra)--(rrt1); \draw[edge] (rra)--(rrb); \draw[edge] (rrb)--(rrt2); \draw[edge] (rrb)--(rrc); \draw[edge] (rrc)--(rrt3); \draw[edge] (rrc)--(rrt4);
\draw[rot] (11.0,7.35)--node[above,font=\scriptsize]{左旋}(12.0,7.35);
\node[n,draw=themecurve] (rr2b) at (13.2,8.5) {$B$}; \node[n] (rr2a) at (12.35,7.4) {$A$}; \node[n] (rr2c) at (14.05,7.4) {$C$};
\node[sub] (rr2t1) at (11.85,6.35) {$T_1$}; \node[sub] (rr2t2) at (12.85,6.35) {$T_2$}; \node[sub] (rr2t3) at (13.55,6.35) {$T_3$}; \node[sub] (rr2t4) at (14.55,6.35) {$T_4$};
\draw[edge] (rr2b)--(rr2a); \draw[edge] (rr2b)--(rr2c); \draw[edge] (rr2a)--(rr2t1); \draw[edge] (rr2a)--(rr2t2); \draw[edge] (rr2c)--(rr2t3); \draw[edge] (rr2c)--(rr2t4);

\draw[themegray,line width=.45pt] (0,5.1)--(15.3,5.1);

% LR
\node[font=\bfseries,text=themeamber] at (3.7,4.55) {LR：先左旋 $B$，再右旋 $A$};
\node[hot] (lra) at (1.8,3.75) {$A$}; \node[n] (lrb) at (1.0,2.8) {$B$}; \node[n] (lrc) at (1.55,1.8) {$C$};
\node[sub] (lrt1) at (.45,1.8) {$T_1$}; \node[sub] (lrt2) at (1.15,.9) {$T_2$}; \node[sub] (lrt3) at (1.95,.9) {$T_3$}; \node[sub] (lrt4) at (2.65,2.8) {$T_4$};
\draw[edge] (lra)--(lrb); \draw[edge] (lra)--(lrt4); \draw[edge] (lrb)--(lrt1); \draw[edge] (lrb)--(lrc); \draw[edge] (lrc)--(lrt2); \draw[edge] (lrc)--(lrt3);
\draw[rot] (3.0,2.5)--node[above,font=\scriptsize]{双旋}(4.0,2.5);
\node[n,draw=themecurve] (lr2c) at (5.2,3.75) {$C$}; \node[n] (lr2b) at (4.35,2.65) {$B$}; \node[n] (lr2a) at (6.05,2.65) {$A$};
\node[sub] (lr2t1) at (3.85,1.6) {$T_1$}; \node[sub] (lr2t2) at (4.85,1.6) {$T_2$}; \node[sub] (lr2t3) at (5.55,1.6) {$T_3$}; \node[sub] (lr2t4) at (6.55,1.6) {$T_4$};
\draw[edge] (lr2c)--(lr2b); \draw[edge] (lr2c)--(lr2a); \draw[edge] (lr2b)--(lr2t1); \draw[edge] (lr2b)--(lr2t2); \draw[edge] (lr2a)--(lr2t3); \draw[edge] (lr2a)--(lr2t4);

% RL
\node[font=\bfseries,text=themeamber] at (11.8,4.55) {RL：先右旋 $B$，再左旋 $A$};
\node[hot] (rla) at (9.0,3.75) {$A$}; \node[n] (rlb) at (9.8,2.8) {$B$}; \node[n] (rlc) at (9.25,1.8) {$C$};
\node[sub] (rlt1) at (8.15,2.8) {$T_1$}; \node[sub] (rlt2) at (8.85,.9) {$T_2$}; \node[sub] (rlt3) at (9.65,.9) {$T_3$}; \node[sub] (rlt4) at (10.35,1.8) {$T_4$};
\draw[edge] (rla)--(rlt1); \draw[edge] (rla)--(rlb); \draw[edge] (rlb)--(rlc); \draw[edge] (rlb)--(rlt4); \draw[edge] (rlc)--(rlt2); \draw[edge] (rlc)--(rlt3);
\draw[rot] (11.0,2.5)--node[above,font=\scriptsize]{双旋}(12.0,2.5);
\node[n,draw=themecurve] (rl2c) at (13.2,3.75) {$C$}; \node[n] (rl2a) at (12.35,2.65) {$A$}; \node[n] (rl2b) at (14.05,2.65) {$B$};
\node[sub] (rl2t1) at (11.85,1.6) {$T_1$}; \node[sub] (rl2t2) at (12.85,1.6) {$T_2$}; \node[sub] (rl2t3) at (13.55,1.6) {$T_3$}; \node[sub] (rl2t4) at (14.55,1.6) {$T_4$};
\draw[edge] (rl2c)--(rl2a); \draw[edge] (rl2c)--(rl2b); \draw[edge] (rl2a)--(rl2t1); \draw[edge] (rl2a)--(rl2t2); \draw[edge] (rl2b)--(rl2t3); \draw[edge] (rl2b)--(rl2t4);

\node[anchor=west,font=\scriptsize,text=themegray] at (0,-.05)
  {图只表达局部指针重组与中序不变量；$T_1\ldots T_4$ 的内部形状与高度未画出。旋转后仍需重新计算高度并检查 $\lvert BF\rvert\le1$。};
\end{tikzpicture}
}
\end{center}
四个面板只改变从失衡点到新增分支的两次方向；$T_1\ldots T_4$ 只表示保持原有中序区间的子树，不承诺它们的具体形状。

AVL 高度界可用最小节点数递推验证。令 $N(h)$ 为高度为 $h$ 的 AVL 树最少节点数，则
\[
N(h)=1+N(h-1)+N(h-2),\qquad N(0)=1,;N(1)=2,
\]
因此 $N(h)$ 至少按 Fibonacci 数增长，反推得到 $h=O(\log n)$。这个递推说明“平衡因子不超过 $1$”确实足以给出对数高度；它不等于每棵 AVL 都是完全树。

红黑树把严格高度平衡放宽为颜色约束：根与空叶为黑；红节点的孩子不能为红；从任一节点到其后代空叶的每条路径具有相同黑高。由此最长根叶路径不超过最短路径两倍，树高仍为 $O(\log n)$。插入冲突按叔节点颜色和结构方向决定变色或旋转；删除若移走黑色会产生黑高亏损，修复目标是把亏损上移、吸收或消除。

红黑树的五条颜色不变量应和“空叶也算黑”一起检查：根黑、外部 NIL 黑、红节点不得有红孩子、任一节点到所有后代 NIL 的黑节点数相同、仍保持 BST 有序。它的旋转比 AVL 少严格度，但通常更新路径更短；不能由“高度是对数级”反推出两者的旋转次数或颜色状态相同。

红黑树插入时，新节点先染红以保持黑高不变；若父节点为黑则结束，若父为红则看叔节点：叔红时父叔变黑、祖父变红并把冲突上移；叔黑时按 LL/LR/RR/RL 方向旋转并重新着色。删除时真正困难的是删除黑节点造成的“额外黑色”：兄弟为红先旋转把情形转成兄弟为黑；兄弟黑且两个孩子黑则把黑色上移；若远侄为黑、近侄为红先对兄弟旋转；远侄为红时借用其颜色完成最终旋转。纸面题应始终检查五条不变量，而不是只背某个 case 的图。

\begin{warnbox}{AVL 与红黑树不能共享旋转口诀}
AVL 的触发量是高度差，红黑树的触发量是颜色与黑高。二者旋转都保持 BST 次序，但合法状态不同；只看到“失衡就旋转”无法判断变色与停止条件。
\end{warnbox}

\section{B 树：用节点扇出降低外存访问高度}
设 $m$ 阶 B 树，每个节点至多有 $m$ 个孩子、至多 $m-1$ 个关键字；除根外的非叶节点至少有 $\lceil m/2\rceil$ 个孩子，所有叶节点位于同一层。节点内关键字有序并划分孩子区间。查找先在节点内确定区间，再沿一个孩子下降；昂贵的是访问多少个节点/块。

“阶”的教材口径可能用最大孩子数 $m$，也可能用最大关键字数；下限、分裂位置和高度计算必须先抄出题目定义。若采用最大孩子数口径，非根节点关键字下限为 $\lceil m/2\rceil-1$；根可少于此数但不能违反空树/单孩子的特殊条件。节点内查找既可顺序比较也可折半，外存复杂度通常只计块访问，不能把节点内比较次数与 I/O 次数混为一项。

插入总在叶层落位。节点溢出时把中间关键字上升父节点，左右关键字分成两个合法节点；父节点若再溢出就继续向上，根分裂会使树高增加 1。删除后若节点低于下限，先向兄弟借一个关键字并经父节点重分配；不能借才与兄弟及父分隔关键字合并。若根最终无关键字且只有一个孩子，用该孩子替代根，树高减 1。

\begin{center}
\resizebox{0.94\linewidth}{!}{\providecolor{themebg}{HTML}{FAFAF7}
\providecolor{themefg}{HTML}{111111}
\providecolor{themegray}{HTML}{666666}
\providecolor{themecurve}{HTML}{1D63B8}
\providecolor{themealert}{HTML}{C53030}
\providecolor{themeamber}{HTML}{B86E00}
\begin{tikzpicture}[>=Stealth,font=\small,color=themefg,text=themefg]
\tikzset{
  keynode/.style={draw=themefg,rounded corners=2pt,fill=themebg,minimum height=7mm,inner xsep=5pt,line width=.8pt},
  over/.style={keynode,draw=themealert,fill=themealert!10!themebg},
  good/.style={keynode,draw=themecurve,fill=themecurve!10!themebg},
  parent/.style={keynode,draw=themeamber,fill=themeamber!10!themebg},
  edge/.style={themegray,line width=.9pt},
  action/.style={->,line width=1.1pt,themegray}
}
\node[anchor=west,font=\bfseries\large] at (0,10.25) {B 树局部修复：溢出就分裂上提，下溢先借位，借不到再合并};
\node[anchor=west,text=themegray] at (0,9.7)
  {以 5 阶 B 树为例：每节点至多 5 个孩子、至多 4 个关键字；非根节点至少 2 个关键字。};

% split panel
\node[anchor=west,font=\bfseries,text=themecurve] at (0,8.75) {1. 插入溢出：5 个关键字超过上限 4};
\node[over] (full) at (4.0,7.75) {$10\mid20\mid30\mid40\mid50$};
\draw[action] (6.55,7.75)--node[above,font=\scriptsize]{中间 key 上提}(8.15,7.75);
\node[parent] (p30) at (10.0,8.45) {$30$};
\node[good] (l20) at (8.9,6.95) {$10\mid20$};
\node[good] (r20) at (11.1,6.95) {$40\mid50$};
\draw[edge] (p30)--(l20); \draw[edge] (p30)--(r20);
\node[anchor=west,font=\scriptsize,text=themegray] at (12.5,7.7)
  {父节点若也溢出，沿同一规则继续向上；根分裂使树高 $+1$。};

\draw[themegray,line width=.45pt] (0,6.1)--(16.7,6.1);

% borrow panel
\node[anchor=west,font=\bfseries,text=themecurve] at (0,5.35) {2. 删除下溢：相邻兄弟有富余时先借位};
\node[parent] (bp1) at (3.5,4.4) {$30$};
\node[keynode] (bl1) at (2.1,3.15) {$10\mid20\mid25$};
\node[over] (br1) at (4.9,3.15) {$40$};
\draw[edge] (bp1)--(bl1); \draw[edge] (bp1)--(br1);
\node[anchor=west,font=\scriptsize,text=themealert] at (5.85,3.15) {右节点仅 1 个 key，低于下限 2};
\draw[action] (7.7,4.05)--node[above,font=\scriptsize]{兄弟最大 key 上父；父分隔 key 下移}(9.7,4.05);
\node[parent] (bp2) at (11.5,4.4) {$25$};
\node[good] (bl2) at (10.1,3.15) {$10\mid20$};
\node[good] (br2) at (12.9,3.15) {$30\mid40$};
\draw[edge] (bp2)--(bl2); \draw[edge] (bp2)--(br2);
\node[anchor=west,font=\scriptsize,text=themegray] at (14.05,3.15) {两侧重新合法，父分隔范围也同步更新};

\draw[themegray,line width=.45pt] (0,2.35)--(16.7,2.35);

% merge panel
\node[anchor=west,font=\bfseries,text=themecurve] at (0,1.65) {3. 兄弟也在下限：借不到，只能把父分隔 key 一起合并};
\node[parent] (mp1) at (3.5,.7) {$30$};
\node[keynode] (ml1) at (2.15,-.45) {$10\mid20$};
\node[over] (mr1) at (4.85,-.45) {$40$};
\draw[edge] (mp1)--(ml1); \draw[edge] (mp1)--(mr1);
\draw[action] (6.55,.25)--node[above,font=\scriptsize]{合并}(8.15,.25);
\node[good] (merged) at (10.3,-.05) {$10\mid20\mid30\mid40$};
\node[anchor=west,font=\scriptsize,text=themegray] at (12.6,-.05)
  {父节点少 1 个 key；若父继续下溢则递归修复。若空父是根，则合并节点成为新根，树高 $-1$。};

\node[anchor=west,font=\scriptsize,text=themegray] at (0,-1.35)
  {图只展示局部节点占用与父分隔 key 的移动；各关键字之间的孩子子树未画出，但其区间顺序必须随 key 重分配一起保持。};
\end{tikzpicture}
}
\end{center}
图只展示局部节点占用与父分隔关键字的移动；各关键字之间的孩子子树没有展开，但修复后仍必须满足孩子区间正确与所有叶同层。

统一不变量是：节点占用合法、节点内有序、孩子区间正确、所有叶同层。分裂和合并是在保持这些不变量下控制高度。

\section{B+ 树：内部节点导航，叶层拥有完整检索序列}
B+ 树把记录或记录指针集中在叶节点；非叶节点只保存导航分隔 key，关键字可能在内部和叶层重复。所有叶节点按 key 顺序链接，因此等值查找沿根到叶完成，范围查找定位起点后沿叶链扫描。

在本题常用的“最大关键字复写”口径中，B+ 树内部节点的关键字数等于子树数，内部索引 key 只是分界副本；即使在内部节点比较相等，也必须继续下降到叶节点取真正记录。叶节点分裂时把分界 key 复制到父节点而不从叶中删除；叶链因此始终保有完整有序序列。删除叶中最大 key 时，父节点副本可继续作为合法分界，不必把它误当成记录。

B 树和 B+ 树查找都为 $O(\log_m n)$，B+ 树的优势主要是更大扇出、固定根到叶路径、$O(\log_m n+k)$ 的范围扫描和沿叶链的全表遍历，而不是渐近阶发生改变。磁盘页预算应把“关键字、指针、记录”分别计数：内部节点不放记录，才能在同一页容纳更多导航项。

\begin{center}\small
\begin{tabularx}{\linewidth}{L{2.0cm}C{.35cm}L{2.0cm}YY}\toprule
概念 A&$\neq$&概念 B&真正判据与题面信号&混淆后果\\\midrule
B 树&$\neq$&B+ 树&B 树各层可命中记录；B+ 树内部只导航、叶链承载全体 key&误判成功层次和范围扫描路径\\
树的阶&$\neq$&节点关键字数&阶限制最大孩子数，关键字上限通常少 1&占用上下限算错\\
分裂&$\neq$&旋转借位&分裂制造新节点；借位在相邻节点和父分隔 key 间重分配&高度变化判断错\\\bottomrule
\end{tabularx}\end{center}

\section{Trie：把前缀共享下沉到路径}
Trie 把字符串的第 $i$ 个字符放在根到深度 $i+1$ 的边上；节点至少包含孩子映射和终止标记 \code{isEnd}。插入、完整词查找与前缀查找都沿字符路径前进，时间为 $O(L)$，其中 $L$ 是查询串长度，而不是词典中的词数。\code{isEnd} 不能省略：存在路径 \code{app}$\to$\code{apple} 不代表短串本身已被插入。

孩子表示决定空间与常数。固定小字符集可用定长数组，转移接近直接定位但每个节点都支付字母表空间；稀疏或 Unicode 字符集可用 Hash/有序映射，只为真实边付费但增加映射成本。总节点数不超过所有已插入字符串长度之和加根节点，前缀共享越多越节省。删除一个词时先清除终止标记，只有节点既非其他词终点又没有孩子时才能向上回收，不能删除共享前缀。

Trie Owns 前缀路径语义，不替代 DS10 的精确等值 Hash，也不替代 B+ 树的有序范围索引。若要自动补全，可定位前缀节点后遍历其子树；若要统计前缀频次，可在路径节点维护经过计数，但插入、删除必须同步维护该派生状态。

\section{代码与测试证据}
完整实现位于 \codepath{code/ds09_ordered_index.hpp}，测试位于 \codepath{code/ds09_ordered_index_test.cpp}。本册用边界明确的 BST 作为可执行核心，平衡树和 B/B+ 的分裂规则保留为机制扩展。
\lstinputlisting[language=C++,firstline=12,lastline=104,firstnumber=1,numbers=left,caption={BST 查找、插入、删除与有序输出}]{30_408/10_数据结构/09_查找与有序索引/code/ds09_ordered_index.hpp}
\lstinputlisting[language=C++,firstline=8,lastline=44,firstnumber=1,numbers=left,caption={空树、重复 key 与删除分支测试}]{30_408/10_数据结构/09_查找与有序索引/code/ds09_ordered_index_test.cpp}
\begin{lstlisting}[language=bash]
clang++ -std=c++17 -Wall -Wextra -Werror -pedantic code/ds09_ordered_index_test.cpp -o /tmp/ds09_test
/tmp/ds09_test
\end{lstlisting}

\section{复杂度与概念边界}
\begin{center}\small
\begin{tabularx}{\linewidth}{L{2.5cm}C{1.8cm}C{1.8cm}YY}\toprule
结构/操作 & 平均或保证 & 最坏 & 依赖的不变量 & 维护代价\\\midrule
顺序查找 & $O(n)$ & $O(n)$ & 无序也可 & 几乎无预处理\\
折半查找 & $O(\log n)$ & $O(\log n)$ & 有序+随机访问 & 插入移动\\
BST & $O(h)$ & $O(n)$ & 局部有序 & 可能退化\\
平衡树 & $O(\log n)$ & $O(\log n)$ & 高度/颜色约束 & 旋转/分裂\\\bottomrule\end{tabularx}\end{center}

\section{做题控制协议}
先写返回契约和区间/树不变量，再看输入是否有序、是否支持随机访问、更新是否频繁。看到 Hash 适合精确等值但不适合范围查询；看到动态有序更新，才考虑平衡索引。每个复杂度都把高度或有序前提写出来。
\begin{corebox}{压缩复原}
五问：查询要返回什么？候选区间怎样缩小？秩序存在哪里？更新破坏哪个不变量？维护成本由谁承担？
\end{corebox}
\end{document}
